\documentclass[12pt]{article}
\usepackage{amsmath, amssymb}
\usepackage{geometry}
\geometry{margin=1in}

\title{Lecture 18 Scribe – Marginal PMF, Correlation, Covariance, Joint Gaussian RVs}
\date{}

\begin{document}

\maketitle

\textit{(Exam Preparation Notes – Step-by-Step Reasoning)}

\hrule

\section*{1. Marginal PMF}

\subsection*{1.1 Definition}

\textbf{The marginal PMF of one variable is obtained by summing the joint PMF over all possible values of the other variable.}

\subsubsection*{Step-by-step reasoning}

\begin{enumerate}
\item Start with a \textbf{joint PMF} $(P(X,Y))$, which gives probabilities of pairs $((x,y))$.
\item Suppose we are interested only in \textbf{one variable}, say $(X)$.
\item The probabilities in the joint PMF include the influence of $(Y)$.
\item To isolate $(X)$, we:
\begin{itemize}
\item \textbf{Fix $(X = x)$}
\item \textbf{Sum over all possible values of $(Y)$}
\end{itemize}
\end{enumerate}

Thus, we remove the dependency on $(Y)$ by aggregation.

\subsection*{1.2 Marginal PMF from Joint PMF Table}

\textbf{Procedure}

Given a table of $(P(X,Y))$:

\subsubsection*{To compute $(P_X(x))$:}
\begin{enumerate}
\item Select a fixed value $(x)$
\item Identify all entries where $(X = x)$
\item Sum all corresponding probabilities across different $(y)$
\end{enumerate}

\subsubsection*{To compute $(P_Y(y))$:}
\begin{enumerate}
\item Fix $(y)$
\item Sum all entries across different $(x)$
\end{enumerate}

\subsection*{1.3 Example (Logical Derivation)}

Given a joint PMF:

\subsubsection*{Solution steps}

\begin{enumerate}
\item List all values of $(X)$ and $(Y)$
\item For each value of $(X)$:
\begin{itemize}
\item Add probabilities across all $(Y)$
\end{itemize}
\item For each value of $(Y)$:
\begin{itemize}
\item Add probabilities across all $(X)$
\end{itemize}
\end{enumerate}

\subsubsection*{Conclusion}

The resulting sums form:
\begin{itemize}
\item Marginal PMF of $(X)$
\item Marginal PMF of $(Y)$
\end{itemize}

\section*{2. Correlation}

\subsection*{2.1 Definition (Product Moment)}

\textbf{A basic way to measure joint behavior is through the product moment.}

It computes:
\begin{itemize}
\item The \textbf{average value of the product $(XY)$}
\end{itemize}

\subsection*{2.2 Interpretation}

\subsubsection*{Step-by-step understanding}

\begin{enumerate}
\item For each outcome:
\begin{itemize}
\item Multiply $(X)$ and $(Y)$
\end{itemize}
\item Take expectation:
\begin{itemize}
\item This gives $(E[XY])$
\end{itemize}
\item This reflects:
\begin{itemize}
\item Whether large values of $(X)$ occur with large values of $(Y)$
\end{itemize}
\end{enumerate}

\subsection*{2.3 Limitations}

The measure:

\[
R_{XY} = E[XY]
\]

\subsubsection*{Issues (step-by-step)}

\begin{enumerate}
\item \textbf{Depends on scale}
\begin{itemize}
\item If $(X)$ or $(Y)$ is scaled, value changes
\end{itemize}
\item \textbf{Computed about origin}
\begin{itemize}
\item Does not consider mean values
\end{itemize}
\item \textbf{Mixes multiple effects}
\begin{itemize}
\item Mean
\item Spread
\item Dependence
\end{itemize}
\end{enumerate}

\subsubsection*{Special case}

\[
R_{XY} = 0
\Rightarrow \text{orthogonal about the origin}
\]

\subsection*{2.4 Transition to Covariance}

To fix these issues:

\begin{enumerate}
\item Remove the effect of mean
\item Center variables:
\begin{itemize}
\item Replace $(X)$ with $(X - E[X])$
\item Replace $(Y)$ with $(Y - E[Y])$
\end{itemize}
\end{enumerate}

This leads to \textbf{covariance}.

\section*{3. Covariance}

\subsection*{3.1 Definition}

\textbf{We first subtract the mean from each variable and then measure how the centered quantities vary together.}

\subsection*{3.2 Step-by-step Derivation}

\begin{enumerate}
\item Compute mean values:
\begin{itemize}
\item $(E[X])$, $(E[Y])$
\end{itemize}
\item Form deviations:
\begin{itemize}
\item $(X - E[X])$
\item $(Y - E[Y])$
\end{itemize}
\item Multiply deviations:
\begin{itemize}
\item $((X - E[X])(Y - E[Y]))$
\end{itemize}
\item Take expectation:
\begin{itemize}
\item This gives covariance
\end{itemize}
\end{enumerate}

\subsection*{3.3 Interpretation}

Covariance captures:
\begin{itemize}
\item \textbf{Positive covariance}:
\begin{itemize}
\item Both variables deviate in same direction
\end{itemize}
\item \textbf{Negative covariance}:
\begin{itemize}
\item Variables deviate in opposite directions
\end{itemize}
\end{itemize}

\subsection*{3.4 Algebraic Form (Step-by-step)}

Start with definition:

\begin{enumerate}
\item Expand product:
\[
(X - E[X])(Y - E[Y])
\]
\item Multiply terms
\item Apply expectation
\item Use linearity of expectation
\item Simplify
\end{enumerate}

This gives a computationally simpler expression.

\subsection*{3.5 Properties of Covariance}

\subsubsection*{1. Symmetry}

\[
\text{Cov}(X,Y) = \text{Cov}(Y,X)
\]

\subsubsection*{2. Variance as Special Case}

\[
\text{Var}(X) = \text{Cov}(X,X)
\]

\subsubsection*{3. Linearity w.r.t Constants}

Scaling variables scales covariance accordingly

\subsubsection*{4. Independence}

If $(X)$ and $(Y)$ are independent:

\[
\text{Cov}(X,Y) = 0
\]

\textbf{Important conclusion}

\begin{itemize}
\item Zero covariance $\Rightarrow$ \textbf{uncorrelated}
\item But \textbf{does NOT imply independence}
\end{itemize}

\subsection*{3.6 Example 1 (Step-by-step)}

\textbf{Given}

Expressions for $(X)$, $(Y)$

\subsubsection*{Solution steps}

\begin{enumerate}
\item Write covariance formula
\item Substitute given expressions
\item Expand terms
\item Apply expectation operator
\item Simplify step-by-step
\end{enumerate}

\subsubsection*{Conclusion}

Final value obtained after systematic substitution

\subsection*{3.7 Example 2 (Continuous Case)}

Given a joint PDF.

\subsubsection*{Step-by-step solution}

\begin{enumerate}
\item Compute:
\begin{itemize}
\item $(E[X])$ using integration
\end{itemize}
\item Compute:
\begin{itemize}
\item $(E[Y])$
\end{itemize}
\item Compute:
\begin{itemize}
\item $(E[XY])$
\end{itemize}
\item Substitute into covariance formula
\item Simplify result
\end{enumerate}

\section*{4. Pearson’s Correlation Coefficient}

\subsection*{4.1 Definition}

\textbf{Pearson’s correlation coefficient calculates the effect of change in one variable when the other variable changes.}

\subsection*{4.2 Step-by-step Interpretation}

\begin{enumerate}
\item Start from covariance
\item Normalize using:
\begin{itemize}
\item Standard deviation of $(X)$
\item Standard deviation of $(Y)$
\end{itemize}
\item This removes unit dependence
\end{enumerate}

\subsection*{4.3 Properties}

Measures:
\begin{itemize}
\item \textbf{Strength} of relationship
\item \textbf{Direction} of relationship
\end{itemize}

It is:
\begin{itemize}
\item \textbf{Unit-free}
\item \textbf{Normalized covariance}
\end{itemize}

\section*{5. Jointly Gaussian Random Variables}

\subsection*{5.1 Definition}

\textbf{A pair $((X,Y))$ is jointly Gaussian if:}

\begin{itemize}
\item The joint density has bivariate Gaussian form
\item Marginals are Gaussian
\end{itemize}

\subsection*{5.2 Intuition}

\subsubsection*{Step-by-step reasoning}

\begin{enumerate}
\item Many real-world variables are influenced by:
\begin{itemize}
\item Multiple small random factors
\end{itemize}
\item When combined:
\begin{itemize}
\item They tend to follow Gaussian behavior
\end{itemize}
\item Hence, joint distribution becomes Gaussian
\end{enumerate}

\subsection*{5.3 Applications}

\begin{itemize}
\item Signal processing
\item Communication systems
\item Noisy measurements
\end{itemize}

\subsection*{5.4 Example: Climate System}

Variables:
\begin{itemize}
\item $(X)$: CO$_2$ concentration
\item $(Y)$: Temperature anomaly
\end{itemize}

\subsubsection*{Reasoning}

\begin{enumerate}
\item Both depend on multiple random processes:
\begin{itemize}
\item Weather
\item Ocean circulation
\item Emissions
\end{itemize}
\item These combined uncertainties:
\begin{itemize}
\item Produce approximately Gaussian behavior
\end{itemize}
\item Hence:
\begin{itemize}
\item Joint distribution $\approx$ Gaussian
\end{itemize}
\end{enumerate}

\subsection*{5.5 Algebraic Form and Interpretation}

Joint Gaussian is characterized by:
\begin{enumerate}
\item Means
\item Variances
\item Covariance
\end{enumerate}

These parameters fully describe:
\begin{itemize}
\item Shape
\item Spread
\item Dependence structure
\end{itemize}

\section*{Final Exam Summary}

\subsection*{Core Flow of Concepts}

\begin{enumerate}
\item Joint PMF $\rightarrow$ Marginal PMF (by summation)
\item Raw correlation $(E[XY]) \rightarrow$ not centered
\item Center variables $\rightarrow$ Covariance
\item Normalize covariance $\rightarrow$ Pearson coefficient
\item Extend to continuous case $\rightarrow$ Joint Gaussian RVs
\end{enumerate}

\subsection*{Key Insights to Remember}

\begin{itemize}
\item Marginalization removes dependence on one variable
\item Covariance measures \textbf{joint deviation from mean}
\item Zero covariance $\neq$ independence
\item Pearson coefficient is \textbf{scale-free}
\item Joint Gaussian models real-world continuous systems
\end{itemize}

\end{document}